\chapter{Tâche 2 — Portage du framework pour le support multi-forks}
\label{chap:t2}
% Entreprise : Free Mobile

\section{Contexte et dette technique initiale}
% TODO : pourquoi plusieurs \gls{fork} du programme initial existent-ils ?
% Quelle dette cela crée-t-il (duplication de tests, divergence, double
% maintenance) ? C'est le « problème business » à résoudre.
\emph{[À rédiger.]}

\section{Objectif de la mission}
% TODO : faire évoluer le framework (tâche 1) pour qu'il supporte
% plusieurs forks avec un socle de tests mutualisé.
\emph{[À rédiger.]}

\section{Conception : factorisation et paramétrage}
% TODO : stratégie de portage (abstraction des différences entre forks,
% configuration par fork, tests communs vs spécifiques). Schéma avant/après.
\emph{[À rédiger.]}

\section{Réalisation}
% TODO : extrait montrant le mécanisme de paramétrisation par fork.
% \begin{lstlisting}[language=Python,caption={Paramétrage d'un test par fork.}]
% ...
% \end{lstlisting}
\emph{[À rédiger.]}

\section{Difficultés rencontrées}
% TODO : divergences de comportement entre forks, rétro-compatibilité.
\emph{[À rédiger.]}

\section{Résultats et bénéfice business}
% TODO : maintenance unifiée, non-régression sur tous les forks d'un coup,
% réduction de la duplication. Relier au thème dette technique.
\emph{[À rédiger.]}
